\section{Related Work}
\label{sec:related}

\subsection{Rail emissions estimation}

Railway greenhouse-gas accounting divides on whether traction fuel is metered. Where an
operator publishes its own fuel consumption, as in most national rail systems, IPCC's Tier~1 is
already close to a measurement: the uncertainty in the estimate is essentially the uncertainty
in the emission factor. Where fuel is not published---typically for fragmented or private
freight networks---inventories model traction energy from gross tonne-kilometres and specific
energy consumption, and the resulting figures carry the uncertainty of a fleet model.

Pakistan sits firmly in the first case, which is why this paper spends no space on traction
modelling and considerable space on what to do when a published national figure disagrees with
the metered one.

\subsection{Inventory verification against national submissions}

Independent verification of national submissions is normally attempted at the level of totals,
using atmospheric inversions or independent global products such as EDGAR
\citep{edgar2025booklet}, and it usually produces a disagreement that cannot be adjudicated: two
models differ, and neither is a measurement. This inventory's own fugitives sector is an
example, where a factor of \num{1.3} against EDGAR settles nothing because both estimates
choose within the same published range.

Railways is different in kind. The disagreement here is against an activity figure that has a
physical bound---the carbon in a quantity of diesel a state operator reports consuming---and a
second, independent bound in the traffic statistics published in the same document. A
submission can be refuted on arithmetic when the arithmetic is this short, and that possibility
is rarer in inventory work than the volume of comparison studies suggests.

\subsection{Emission-factor tables as a failure mode}

The literature on inventory quality concentrates on activity data and on factor values. It
says much less about factor \emph{selection}---which of several published tables applies to a
given category---even though the guidelines themselves contain several near-duplicate tables
that differ only in their non-\COtwo{} rows.

This inventory has now documented three instances: residential fuel against the power table,
power against several, and rail against road here. The common structure is that the \COtwo{}
row is identical across tables because \COtwo{} depends only on the carbon in the fuel, so the
dominant number in the category is unaffected and every obvious check passes. Section
\ref{sec:factors} treats the rail case and Section~\ref{sec:discussion} the general one.
